\documentclass{article}

\input{../definitions.tex}

\usepackage{mathpartir}
\usepackage{fvextra}
\newdef{problem}
\usepackage[a4paper]{geometry}
\newpf{solution}{}
%\usepackage[a4paper]{geometry}
%\NewDocumentEnvironment{solution}{+b}{}{}

\RequirePackage{mathtools}
\RequirePackage{unicode-math}
\RequirePackage{polyglossia}

\setmainfont{Latin Modern Roman}
% \setmainfont{TeX Gyre Pagella}
% \setmathfont{TeX Gyre Pagella Math}
\setmathfont{Latin Modern Math}
\setmathfont{Asana Math}[range={scr}]
\setmathfont{STIX Two Math-Regular}[range={bb,"1D538-"1D56B,"0211D}]
\setmathfont{STIX Two Math-Regular}[range={"2102,"210D,"2115,"2119,"211A,"211D,"2124}]
% \setmathfont{TeX Gyre Pagella Math}[range={8730}] %sqrt
% \setmathfont{Asana Math}[range={"007B,"007D}]  % {}
\setmathfont{Asana Math}[range={8709,\setminus,"2216,"29F5,"219D}]  % U+2205, emptyset

\newcommand{\seq}[3]{#1 ⊢ #2 : #3}
\def\emp{\bullet}
\newcommand{\lam}[3]{λ #1 : #2.\: #3}
\renewcommand{\pair}[2]{\p{#1, #2}}
\newcommand{\match}[5]{\textsf{match}\:{#1}\:| \inl {#2} ↦ {#3} | \inr {#4} ↦ {#5}}
\newcommand{\emptybox}[1]{\boxed{\phantom{#1}}}
\newcommand{\nullaryrule}[2]{
  \ExpandArgs{c}\DeclareDocumentCommand{#1}{m}{\inferrule*[right=\rlap{\(#2\)}] { } {##1}}
}
\newcommand{\unaryrule}[2]{
  \ExpandArgs{c}\DeclareDocumentCommand{#1}{mm}{\inferrule*[right=\rlap{\(#2\)}] {##2} {##1}}
}
\newcommand{\binaryrule}[2]{
  \ExpandArgs{c}\DeclareDocumentCommand{#1}{mmm}{\inferrule*[right=\rlap{\(#2\)}] {##2 \\ ##3} {##1}}
}
\newcommand{\ternaryrule}[2]{
  \ExpandArgs{c}\DeclareDocumentCommand{#1}{mmmm}{\inferrule*[right=\rlap{\(#2\)}] {##2 \\ ##3 \\ ##4} {##1}}
}
%\newcommand{\ax}[1]{\inferrule*[right=\textsc{Ax}] { } {#1}}
\nullaryrule{ax}{\textsc{Ax}}
\binaryrule{prodI}{×\textsc{I}}
\unaryrule{prodEL}{×\textsc{E}₁}
\unaryrule{prodER}{×\textsc{E}₂}
\unaryrule{funI}{→\textsc{I}}
\binaryrule{funE}{→\textsc{E}}
\unaryrule{sumIL}{+\textsc{I}{}₁}
\unaryrule{sumIR}{+\textsc{I}{}₂}
\ternaryrule{sumE}{+\textsc{E}}

\def\formula{\texttt{Formula}}
\def\nnf{\texttt{NNF}}
\def\cnf{\texttt{CNF}}
\def\var{\texttt{Var }}
%\def\til{\sim}
\def\tonnf{\texttt{to-nnf}}
\def\eval{\texttt{eval}}
\def\evalnnf{\texttt{eval-nnf}}
\def\evalcnf{\texttt{eval-cnf}}
\def\assignment{\texttt{Assignment}}
\def\bool{\texttt{Bool}}
\def\maybe{\texttt{Maybe }}
\def\fin{\texttt{Fin }}
\def\lit{\texttt{Literal}}
\def\disj{\texttt{Disjunct}}

\newcommand{\pr}[2]{\textsf{pr}_{#1} {#2}}

\title{Logika v Računalništvu project}
\date{Dated 2026-04-29}
\begin{document}
\maketitle


\begin{problem}[*]
  Define a type of formulas called \formula, with the following grammar:
  \begin{align*}
    &&&&&&&&&&&&&\formula &→&&& \var n&&&&&&&&&&&&\\
    &&&&&&&&&&&&&         &|&&& \neg \formula&&&&&&&&&&&&\\
    &&&&&&&&&&&&&         &|&&& \formula ∧ \formula&&&&&&&&&&&&\\
    &&&&&&&&&&&&&         &|&&& \formula ∨ \formula&&&&&&&&&&&&
  \end{align*}
\end{problem}


\begin{problem}[*]
  Define a type of negation normal form formulas called \nnf, with the following grammar:
  \begin{align*}
    &&&&&&&&&&&\lit &→&&& \var n&&&&&&&&&&&&\\
    &&&&&&&&&&&     &|&&& ¬ \var n&&&&&&&&&&&&\\
    &&&&&&&&&&&\nnf &→&&& \lit&&&&&&&&&&&&\\
    &&&&&&&&&&&     &|&&& \nnf ∧ \nnf&&&&&&&&&&&&\\
    &&&&&&&&&&&     &|&&& \nnf ∨ \nnf&&&&&&&&&&&&
  \end{align*}
\end{problem}

\begin{problem}[*]
  Construct a function \tonnf{} of type \(\formula → \nnf\) that converts a formula to an equivalent formula in negation normal form.

  Note: You may find a more suitable name for the functions.
\end{problem}


\begin{problem}[**]
  Copy the \texttt{Assoc} module from week 9 exercises and complete it to a fully working implementation of an associative structure you want (associative list, dictionary, etc.). For example, use the \texttt{NoDup} structure from the exercises.

  Then use
\begin{Verbatim}[breaklines, breakanywhere, mathescape, codes={\catcode`_=12}]
open Assoc $ℕ$ test-$≡$ Bool
Assignment : Set
Assignment = Assoc
\end{Verbatim}

  Note: it is also OK to specialise the definitions to \(K = ℕ\) and \(V = \bool\), for example to use \(\fin n\).
  You are also free to use the \texttt{stdlib} definitions like
\begin{Verbatim}[breaklines, breakanywhere, mathescape, codes={\catcode`_=12}]
open import Relation.Binary using (Decidable; DecidableEquality)
open import Data.List.Relation.Unary.Any using (Any; any?)
open import Data.List.Relation.Unary.All using (All; all?)
open import Relation.Nullary using (yes; no; ¬_; ¬?)
\end{Verbatim}
\end{problem}

\begin{problem}[*]
  Define an evaluation function \(\eval : \assignment → \formula → \maybe \bool\) assigning to each assignment of variables and formula its truth value.
\end{problem}

\begin{problem}[*]
  Define an evaluation function \(\evalnnf : \assignment → \nnf → \maybe \bool\) assigning to each assignment of variables and negation normal from formula its truth value.
\end{problem}

\begin{problem}[*]
  Define a type of conjunction normal form formulas called \cnf, with the following grammar:
  \begin{align*}
    &&&&&&&&&&&\lit  &→&&& \var n&&&&&&&&&&&&\\
    &&&&&&&&&&&      &|&&& ¬ \var n&&&&&&&&&&&&\\
    &&&&&&&&&&&\disj &→&&& \lit&&&&&&&&&&&&\\
    &&&&&&&&&&&      &|&&& \lit ∨ \disj&&&&&&&&&&&&\\
    &&&&&&&&&&&\cnf  &→&&& \disj ∨ \cnf&&&&&&&&&&&&
  \end{align*}
\end{problem}

\begin{problem}[*]
  Define an evaluation function \(\evalcnf : \assignment → \cnf → \maybe \bool\) assigning to each assignment of variables and conjunction normal from formula its truth value.
\end{problem}

\begin{problem}[**/***]
  Write an SAT solver for \cnf{} formulas. It should output either an assignment such that evaluating the formula at that assignment evaluates to \texttt{true} or that no such assignment exists.

  Note: a more complex implementation (e.~g. DPLL) will be graded higher.
\end{problem}

\begin{problem}[**]
  Show that the SAT solver you implemented is indeed correct, if that is not obvious from the output type of the SAT solver.
\end{problem}

\begin{problem}[**/***]
  Write a function that converts an \nnf{} formula to an equisatisfiable \cnf formula.

  Note: it is intended for you to attempt to implement the Tseytin transformation. A simpler implementation will be accepted for partial credit.
\end{problem}

\begin{problem}[****?]
  Show that the formula produced by the above function is equisatisfiable to the given formula.
\end{problem}



\end{document}


%%% Local Variables:
%%% mode: latex
%%% TeX-master: t
%%% End:
